\documentclass{beamer}

\usepackage{booktabs}
\usepackage{array}

\title{Formal Modeling of the Matter Spec:\\Compositional State-Space Decomposition}
\subtitle{Progress Update}
\author{}
\date{\today}

\begin{document}

\frame{\titlepage}

% ------------------------------------------------------------
\begin{frame}{Context \& Goal}
\begin{itemize}
  \item Pipeline: Matter spec PDF $\rightarrow$ structured corpus $\rightarrow$
        91 ranked candidate sections $\rightarrow$ TLA+ models $\rightarrow$ TLC checking
  \item Existing reference artifact: \texttt{ArmFailSafe.tla}
        (\texttt{core\_11.10.7.2}), hand-built, single section
  \item \textbf{Today's goal:} model more sections without TLC running out
        of memory as models get larger
\end{itemize}
\end{frame}

% ------------------------------------------------------------
\begin{frame}{The Proposed Method — and Why It Was Unsound}
\begin{itemize}
  \item A guide proposed adapting a real paper's decomposition method
        (WCDCAnalyzer, NDSS 2026) to avoid state-space explosion
  \item Citations checked out as real, but the \emph{technique} did not
        match the paper:
  \begin{itemize}
    \item Guide: fix one ``representative value'' at a phase boundary
          (e.g.\ one expiry value out of many possible ones)
    \item Paper: treats handed-off values symbolically/exhaustively,
          under a precondition (disjoint cryptographic primitives) that
          does not even apply to our non-cryptographic models
  \end{itemize}
  \item \textbf{Conclusion:} the guide's shortcut is unsound sampling —
        a bug at an untested value would simply never be found
\end{itemize}
\end{frame}

% ------------------------------------------------------------
\begin{frame}{Our Corrected Method}
\begin{itemize}
  \item Any handed-off variable must range over its \textbf{full domain}
        at a phase boundary — never a fixed sample
  \item Scope: model the 91 already-ranked candidate sections, not the
        full 1335-page spec, not arbitrary hand-picked clusters
  \item Use the corpus's existing cross-reference tool
        (\texttt{DependencyGraph.closure()}) to compose genuinely related
        sections together, instead of treating each in isolation
\end{itemize}
\end{frame}

% ------------------------------------------------------------
\begin{frame}{Pilot Section: Problems Faced \& Fixed}
\texttt{core\_4.20.3.8} — Packet Acknowledgements
\vspace{0.3cm}
\begin{center}
\small
\begin{tabular}{p{4.6cm}p{5.6cm}}
\toprule
\textbf{Problem} & \textbf{Fix} \\
\midrule
Ack outcome split into 3 separate injectable events & Merged into one event; outcome derived from state \\
\addlinespace
Helper action reachable at an impossible time & Removed from the event queue entirely \\
\addlinespace
Automatic reaction modeled as an external event & Folded into the action that actually triggers it \\
\bottomrule
\end{tabular}
\end{center}
\vspace{0.3cm}
Result: 9 reachable states; one genuine finding survived (stale timer
notification) after all self-inflicted artifacts were removed.
\end{frame}

% ------------------------------------------------------------
\begin{frame}{Confirming the Pattern \& Scaling Up}
\begin{itemize}
  \item \texttt{core\_4.19.4.8} (Bluetooth transport) turned out to share
        the exact same design as \texttt{core\_4.20.3.8} (Wi-Fi transport)
  \item Reused the validated model $\rightarrow$ same finding recurs
        identically across two transport bindings
  \item \textbf{Scaled up:} composed 8 real, cross-referenced subsections
        of \texttt{core\_4.20.3} (handshake, ack/window, message
        segmentation, shutdown) into one integrated session model
\end{itemize}
\end{frame}

% ------------------------------------------------------------
\begin{frame}{Results: Decomposition Works}
\begin{center}
\begin{tabular}{lr}
\toprule
\textbf{Model} & \textbf{Reachable states} \\
\midrule
Naive, combined & 4{,}217 \\
\midrule
Decomposed: Handshake & 56 \\
Decomposed: Ack/Window & 348 \\
Decomposed: Segmentation & 152 \\
\bottomrule
\end{tabular}
\end{center}
\vspace{0.3cm}
\begin{itemize}
  \item Peak state space: 4{,}217 $\rightarrow$ 348 (\textbf{$\sim$12x smaller})
  \item Re-checked all three pieces independently: confirmed the same
        genuine findings still surface — decomposition lost nothing
\end{itemize}
\end{frame}

% ------------------------------------------------------------
\begin{frame}{Findings So Far}
\begin{center}
\small
\begin{tabular}{p{1.4cm}p{9.2cm}}
\toprule
\textbf{ID} & \textbf{Finding} \\
\midrule
A & Stale ``timer expired'' notification arriving after the timer was
already stopped — recurs identically in two transport bindings \\
\addlinespace
B & Data/ack events arriving before the handshake completes \\
\addlinespace
C & Stale handshake response/timeout arriving after the session is
already established \\
\addlinespace
D & Connection shutdown described in only 28 words — no procedure at all \\
\bottomrule
\end{tabular}
\end{center}
All are \textbf{candidates for human review}, not confirmed conclusions.
\end{frame}

% ------------------------------------------------------------
\begin{frame}{What to Check Next}
\begin{itemize}
  \item Confirm the source spec's referenced diagrams don't already
        resolve findings A--C (our pipeline can't read figures, only text)
  \item Add the second required catch-all state (contradictory/overlapping
        transitions), not just the ``undefined'' one built so far
  \item Systematically confirm decomposition never silently drops a
        finding, beyond spot-checking known cases
  \item Check whether any finding depends on an invented bound
        (e.g.\ window/queue size) rather than the spec itself
\end{itemize}
\end{frame}

% ------------------------------------------------------------
\begin{frame}{Next Steps, Per Methodology}
\begin{itemize}
  \item Continue through the remaining 88 ranked candidate sections
  \item Cluster sections by real cross-references \emph{before} modeling —
        never compose sections that don't actually interact
  \item Apply sequential (phase-based) decomposition where a flow is
        genuinely one-way (e.g.\ the commissioning flow)
  \item Longer term: multiple independent models per section, treating
        disagreement between them as an ambiguity signal
\end{itemize}
\end{frame}

% ------------------------------------------------------------
\begin{frame}{Summary}
\begin{itemize}
  \item Identified and corrected an unsound shortcut in the proposed
        decomposition method
  \item Validated a sound, lossless alternative on a real composed example
        (8 cross-referenced spec subsections)
  \item Demonstrated a $\sim$12x reduction in peak state space with no
        findings lost
  \item 4 candidate underspecification findings recorded, ready for review
  \item Clear, methodology-aligned roadmap for scaling to the rest of the
        corpus
\end{itemize}
\end{frame}

\end{document}
